\documentclass[journal=jacsat,manuscript=article]{achemso}

\usepackage{amsmath}
\usepackage{amsfonts}
\usepackage{hyperref}
\usepackage{braket}
\usepackage{url}
\usepackage{graphicx}

\DeclareUnicodeCharacter{2009}{\ }

\author{Kyle A. Hamer}
\email{khamer3@lsu.edu}
\affiliation[Louisiana State University]
{Department of Physics and Astronomy, Louisiana State University, Baton Rouge, LA 70803}

\author{Aderonke S. Folorunso}
\affiliation[Louisiana State University]
{Department of Chemistry, Louisiana State University, Baton Rouge, LA 70803}
%0000-0003-4124-3052

\author{Kenneth Lopata}
\affiliation[Louisiana State University]
{Department of Chemistry, Louisiana State University, Baton Rouge, LA 70803}
\altaffiliation{Center for Computation and Technology, Louisiana State University, Baton Rouge, LA 70803}
% 0000-0002-9141-684X

\author{Kenneth J. Schafer}
\affiliation[Louisiana State University]
{Department of Physics and Astronomy, Louisiana State University, Baton Rouge, LA 70803}
% 0000-0003-4740-9113

\author{Mette B. Gaarde}
\email{mgaarde1@lsu.edu}
\affiliation[Louisiana State University]
{Department of Physics and Astronomy, Louisiana State University, Baton Rouge, LA 70803}
% XXX

\author{Fran\c{c}ois Mauger}
\affiliation[Louisiana State University]
{Department of Physics and Astronomy, Louisiana State University, Baton Rouge, LA 70803}
%0000-0001-7555-6001

\title{Tracking Charge Migration with Frequency-Matched Strobo-Spectroscopy}

\begin{document}

% %%%%%%%%%%%%%%%%%%%%%%%%%%%%%%%%%%%%%%%%%%%%%%%%%%%%%%%%% %
% ABSTRACT
% %%%%%%%%%%%%%%%%%%%%%%%%%%%%%%%%%%%%%%%%%%%%%%%%%%%%%%%%% %

\newpage
\begin{abstract}
    We present frequency-matched strobo-spectroscopy (FMSS) of charge migration (CM) in bromodiacetylene, simulated with time-dependent density-functional theory. CM+FMSS is a pump-probe scheme that uses a frequency-matched HHG-driving laser as an independent probe step following the creation of a localized hole on the bromine atom that induces CM dynamics. We show that the delay-dependent harmonic yield tracks the phase of the CM dynamics through its sensitivity to the amount of electron density on the bromine end of the molecule. FMSS takes advantage of the intrinsic attosecond time resolution of the HHG process, in which different harmonics are emitted at different times and thus probe different locations of the electron hole. Finally, we show that the CM-induced modulation of the HHG signal is dominated by the recombination step of the HHG process, with negligible contribution from the ionization step.
\end{abstract}

% %%%%%%%%%%%%%%%%%%%%%%%%%%%%%%%%%%%%%%%%%%%%%%%%%%%%%%%%% %
% INTRODUCTION
% %%%%%%%%%%%%%%%%%%%%%%%%%%%%%%%%%%%%%%%%%%%%%%%%%%%%%%%%% %

\newpage
\section{Introduction}

\noindent Understanding the ultrafast motion of electrons within matter is of critical importance for many areas of science and technology. One example is charge migration \cite{calegari2016, worner2017, nisoli2017, li2020} (CM): the coherent motion of a positively-charged electron hole along the backbone of a molecule following a localized ionization event, which can be observed on Angstrom spatial scales and attosecond time scales \cite{goulielmakis2010, calegari2014, kraus2015}. CM is a widely-studied phenomenon due to its potential for understanding and perhaps steering downstream processes such as chemical reactions, photosynthesis, and photovoltaics via charge-directed reactivity \cite{remacle1998, golubev2015, mauger2022}. Since its discovery in the late 1990s \cite{cederbaum1999, breidbach2003}, the study of CM has flourished \cite{lunnemann2008, golubev2015, kuleff2016, bruner2017, folorunso2021, mansson2021}, with much of this research performed in recent years \cite{barillot2022, li2022, schlegel2022, yong2022, zhao2022, folorunso2023, yu2023, he2022, hamer2022, matselyukh2022, kobayashi2022}. 

Despite the many challenges of doing experiments on the attosecond timescale, CM has been measured using several different techniques, including X-ray absorption spectroscopy \cite{bressler2010, goulielmakis2010, kobayashi2022, matselyukh2022}, photoelectron spectroscopy \cite{bandrauk2004, stolow2004, calegari2014, ayuso2017, lara-astiaso2018}, and high-harmonic spectroscopy (HHS) \cite{bredtmann2012, worner2013, kraus2015, peng2019, he2022}. Due to its inherent sub-femtosecond temporal resolution via the attochirp \cite{dudovich2006, smirnova2009, azoury2017} of the harmonic radiation, in which different harmonic energies are emitted at different times during the laser cycle, HHS is particularly well-suited to perform time-resolved measurements of ultrafast electron dynamics via a pump-probe scheme. It is useful make a distinction between schemes where the CM is initiated by the same laser field that probes those dynamics, in which the CM dynamics is re-initiated every half-laser cycle~\cite{kraus2015, he2022}, and schemes where the pump and probe steps are independent \cite{hamer2022}, as discussed here.

\begin{figure*}[!t]
    \centering
    \includegraphics[width=0.65\linewidth]{Figures/figure1.png}
    \caption{Schematic of our CM+FMSS analysis: (a) snapshots of the time-dependent CM dynamics following the creation of a localized hole on the bromine atom. (b) These dynamics are probed by a delayed, frequency-matched HHG-driving laser field (red curve). The semiclassical return energy of the rescattered electron wave packet is plotted as a function of the recombination time (blue curve). (c) The resulting orbital-resolved, normalized CM+FMSS spectrum (see text). There is a clear variation in the delay-dependent harmonic spectrum due to the CM dynamics.}
    \label{fig:1}
\end{figure*}

In this paper, we present frequency-matched high-harmonic strobo-spectroscopy of charge migration (CM+FMSS), simulated with time-dependent density-functional theory (TDDFT) \cite{kohn1965, runge1984}. We induce CM dynamics in a bromodiacetylene (BrC$_{4}$H) molecule via the creation of a localized hole on the bromine end of the molecule \cite{folorunso2021, hamer2022}. Following the initiation of the CM dynamics, CM+FMSS uses a delay-dependent, few-cycle HHG-driving laser pulse as an independent probe step to precisely determine the time-dependent location of the electron hole, by tracking the amount of electron density on the bromine atom. The driving laser field is polarized perpendicular to the CM motion, so that it does not drive the electron density. We match the frequency $\omega_{L}$ of the laser to $\omega_{\text{CM}}$ such that the position of the electron hole is the same in each half-cycle of the laser field for any given delay. If not already known, we showed in Ref.~\cite{hamer2022} that the CM frequency can be obtained using high-harmonic sideband spectroscopy. 

In Figure~{\ref{fig:1}}, we show a schematic which describes the FMSS concept: panel (a) depicts the time-dependent CM dynamics. This dynamics is probed after some time $\tau$ by a HHG-driving laser field, shown in (b), with $\omega_{L} = \omega_{\text{CM}}/2$. Also in panel (b) we show the semiclassical \cite{corkum1993, schafer1993, lewenstein1994, chang2011} time-dependent return energies during one-half cycle of the HHG-driving laser field, which maps to the sub-cycle time-dependent emission frequencies in the harmonic spectrum \cite{varju2005, niikura2005, smirnova2009}. For different delays, a given harmonic energy is emitted at a different time during the CM period, and this strongly affects the resulting HHG yield as shown in panel (c). For example, at a delay of 4.55 optical cycles (o.c.), the low-order (high-order) harmonics are emitted when the hole is located on the bromine atom (terminal $\text{C}\equiv\text{C}$ bond), which gives rise to low (high) HHG yield -- see the purple dashed lines in Fig.~{\ref{fig:1}}. A central finding of this work is that the delay-dependent harmonic yield tracks the time-dependent electron density on the bromine atom, from which we determine the phase of the CM motion.

\section{Methods}

\begin{figure*}[!t]
    \centering
    \includegraphics[width=0.65\linewidth]{Figures/figure2.png}
    \caption{(a) Schematic of the FMSS configuration to probe CM in BrC$_{4}$H. The HHG driving field is polarized perpendicular to the molecular backbone along which the periodic CM occurs. We also show the isosurface of the initial CM orbital density.
    %
    (b,c) Time evolution of the (b) CM orbital density and (c) electron hole density, integrated over the directions transverse to the molecular backbone, as a function of position.}
    \label{fig:2}
\end{figure*}

\noindent In our CM+FMSS simulations, we start by creating a one-electron valence hole localized on the halogen end of a bromodiacetylene ($\text{BrC}_{4}\text{H}$) molecule. We use constrained density functional theory to create an outer-valence hole at $t=0$ that induces particle-like CM along the backbone of the molecule, with a fundamental frequency of $\omega_{\text{CM}} = 1.85\  \text{eV}$, similar to Refs.~\cite{folorunso2021, hamer2022}. This initial hole emulates a halogen-localized ionization \cite{sandor2019, barillot2022}, which leads to CM in the valence shell of the molecule, as we have demonstrated previously~\cite{bruner2017, folorunso2021, mauger2022, hamer2022}. The CM dynamics is illustrated in Fig.~{\ref{fig:2}}: in panel (a) we show the isosurface of the electron density contribution from the unpaired Kohn-Sham channel from which we remove one electron, here called the \emph{CM orbital} $\psi_{\text{CM}}$. In panels (b) and (c), we show two different representations of the resulting CM dynamics. Fig.~{\ref{fig:2}}(b) shows the time evolution of the CM orbital density, denoted $|\psi_{\text{CM}}(z, t)|^{2}$, integrated over the directions transverse to the molecular backbone. Here, we see a clear oscillation of the electron density in the CM orbital which begins on the bromine atom, travels through the central carbon bond to the terminal bond, and then back again, with a period of $2.24\ \text{fs}$. In Fig.~{\ref{fig:2}}(c), we show the corresponding hole density, defined as the time-dependent density difference between the neutral and the cation densities, $\rho_{h}(z, t) = \rho_{\circ}(z) - \rho_{+}(z, t)$ \cite{cederbaum1999, breidbach2003, folorunso2021}, again integrated over the directions transverse to the molecular backbone. This hole density exhibits a similar pattern to that of the electron density in the CM orbital with additional high-frequency oscillations localized around each of the atomic centers \cite{mauger2022}. In all calculations shown in this paper, nuclear dynamics have been omitted~\cite{footnote1}.

We next induce HHG in the BrC$_{4}$H cation undergoing CM, using a laser pulse with a polarization direction perpendicular to the molecular backbone so that we do not drive the electron density along the molecular backbone. This laser pulse has a frequency $\omega_{L} = \omega_{\text{CM}}/2$ (corresponding to a laser wavelength $\lambda_{L} = 1344\ \text{nm}$). The frequency-matching condition is chosen such that the electron hole is at the same position along the molecular backbone at every half-cycle of the laser field. By using different sub-cycle delays between the initiation of the CM and the laser field, we therefore sample different positions of the electron hole along the molecular backbone. For our TDDFT simulations, we use $\sin^{2}$ laser pulses centered around a delay $\tau$ relative to the initiation of the CM, and that last for 5 o.c. in total ($\approx 1.5$ o.c. FWHM). We then scan the sub-cycle-resolved delay over two full laser cycles, advancing the delay $\tau$ in increments of $1/16$ optical cycles. In all simulations we use a peak intensity of $45\ \text{TW}/\text{cm}^{2}$, leading to a cutoff energy of around $40\ \text{eV}$.

We use grid-based TDDFT with a local-density-approximation exchange-correlation functional \cite{perdew1981, yabana1996, marques2012} and average-density self-interaction correction \cite{legrand2002, ciofini2003, tsuneda2014} within the OCTOPUS software package \cite{marques2003, castro2006, andrade2015, tancogne-dejean2020} to describe both the CM and HHG processes. We use a simulation box with dimensions of $90 \times 40 \times 90\  \text{a.u.}$ (with the shorter box length transverse to both the laser field and molecular axes), and a complex absorbing potential that extends $15\  \text{a.u.}$ from each edge of the box. We choose the box dimensions such that we select the short-trajectory contribution to the HHG spectrum that is usually observed in HHG measurements \cite{hamer2021}, by absorbing the long-trajectory contribution. We use a grid spacing of $0.3\  \text{a.u.}$ in all directions. 

In order to compute harmonic spectra, we first define the orbital-resolved dipole moment \cite{hamer2021} $d_{j}(t)$ corresponding to the $j^{\text{th}}$ Kohn-Sham orbital $\psi_{j}(\vec{r}, t)$:
\begin{equation}\label{eq:ordm}
    \vec{d}_{j}[\tau](t) \cdot \hat{x} = \int x\, \mathrm{d}x \iint \mathrm{d}y\, \mathrm{d}z\, |\psi_{j}[\tau](\vec{r}, t)|^{2}
\end{equation}
Here, we focus on the dipole signal parallel to the driving laser field (in the $x$-direction). We have checked that our results are nearly identical when including the dipole signal in the directions perpendicular to the laser field. The oscillating charge density along the molecular axis induces a significant dipole contribution along the axis of the molecular backbone (the $z$-direction); however, it manifests only in the low-frequency ($\leq 10\  \text{eV}$) regime. In the remainder of this paper, we focus on either the combined dipole signal from the three $\pi$ orbitals in which the CM takes place (and where the vast majority of the harmonic signal resides), or the dipole signal from only the CM orbital defined above. We window the thus-computed dipole moment in the time domain using a $\cos^{2}$ function which has the same width as the laser pulse, such that the time-dependent signal smoothly goes to zero on both ends. Then, we apply a Fourier transform, and square to obtain the delay-dependent HHG yield:
\begin{equation}\label{eq:fmss_spec}
    S[\tau](\omega) = \bigg{|} \mathcal{F} \left[W(t, \tau) \cdot \left(\vec{d}[\tau](t) \cdot \hat{x}\right)\right] \bigg{|}^{2}
\end{equation}
Lastly, in order to more clearly investigate the CM-induced delay-dependent modulation of the harmonic signal, we first smooth the spectrum using a moving average to remove the individual harmonic peaks and then normalize $S[\tau](\omega)$ by the delay-averaged harmonic signal; this final step removes the general shape (perturbative region, plateau, and cutoff region) of the harmonic spectrum and focuses on the delay dependence.

\section{Results and Discussion}

\begin{figure}[!t]
    \centering
    \includegraphics[width=0.65\linewidth]{Figures/figure3.png}
    \caption{(a) Normalized, CM-orbital-resolved FMSS spectrum for $\lambda_{L} = 1344\  \text{nm}$ ($\omega_{L} = \omega_{\text{CM}}/2$), $I_{\circ} = 45\  \text{TW}/\text{cm}^{2}$. (b) Delay-dependent harmonic spectrum taken from model calculations (see text). Black dashed lines are taken from ridge detection of the peaks in panel (b).}
    \label{fig:3}
\end{figure}

\noindent In Figure~\ref{fig:3}(a), we show the normalized CM+FMSS spectrum calculated from the CM-orbital-resolved dipole moment described by Eq.~{\ref{eq:ordm}}, also shown previously in Fig.~{\ref{fig:1}}(c), around 4 optical cycles (approximately $18\, \text{fs}$) after the initiation of the CM. Clearly, there is a pronounced half-laser-cycle-periodic, delay- and harmonic-frequency-dependent variation in the harmonic signal which is not present in the absence of the CM dynamics (in the neutral molecule). This variation is such that the yield is roughly five times more intense when the hole is \emph{not} on the bromine atom. Below the cutoff energy $E_{c} = 40\  \text{eV}$, this spectral maximum trends towards earlier delays as the harmonic frequency increases. As we discuss below, the slope of this tilt matches the negative of the attochirp of the harmonic radiation.

To further investigate our TDDFT results, we construct a CM-modulated model HHG dipole moment based on the strong-field approximation (SFA) \cite{lewenstein1994}. In the absence of CM, the idealized harmonic response from a gas-phase target irradiated by a monochromatic laser field with a frequency $\omega_{L}$, a cutoff frequency $\omega_{c}$, and an envelope $F(t)$ is given by
\begin{equation}\label{eq:dip_noCM}
    d(t) = F(t) \cdot \sum_{n=1,3,5\ldots}^{\infty} A_{n} \sin{\left(n\omega_{L}t + \phi_{n}\right)}
\end{equation}
where the amplitude and phase of the $n^{\text{th}}$ harmonic are respectively defined by
\begin{align}
    A_{n} &= \begin{cases} 1, & n \omega_{L} \leq \omega_{c} \\ e^{-(n\omega_{L} - \omega_{c})}, & \text{otherwise} \end{cases} \label{eq:model_amp}\\
    \phi_{n} &= -\frac{3 \omega_{L}}{10 U_{p}}n^{2} - n \omega_{L} \tau \label{eq:model_phase}
\end{align}
Here, $U_{p}$ is the ponderomotive energy and $\tau$ is the time delay defined above. The harmonic phases $\phi_{n}$ are defined such that we replicate a linear approximation of the semiclassical short-trajectory attochirp \cite{chang2011}. To model the effect of the CM dynamics, we modulate the signal in Eq.~{\ref{eq:dip_noCM}} via:
\begin{equation}\label{eq:dip_CM}
    \tilde{d}(t) = d(t) \cdot \left[1 + \sum_{m} B_{m} \sin{(m \omega_{\text{CM}} t + \phi_{m})}\right]
\end{equation}
where the parameters $\{m, B_{m}, \phi_{m}\}$ describe the individual Fourier components of the field-free CM dynamics.  Consistent with our previous results \cite{hamer2022}, we include two main Fourier components in the CM dynamics: one at $\omega = 1 \cdot \omega_{\text{CM}}$, and a second at $2 \cdot \omega_{\text{CM}}$ which is roughly four times less intense than the first, and with an extra $\pi/4$ phase shift (see again Fig.~{\ref{fig:2}}). 

We plot the delay-dependent harmonic spectrum calculated from the model dipole signal of Eq.~{\ref{eq:dip_CM}} in Figure~{\ref{fig:3}}(b). Like in panel (a), we see a half-cycle-periodic modulation tilting to the left as the harmonic frequency increases. The modulations seen in both panels are consistent with one another, as evidenced by the black dashed lines in both plots, taken from a ridge detection of the peaks in the model spectrum in (b). Removing the attochirp from our model calculations (first term in Eq.~{\ref{eq:model_phase}}) eliminates the slope of the variation shown in Fig.~{\ref{fig:3}}(b). The delay dependence of the variation in the harmonic signal in Fig.~{\ref{fig:3}} is therefore sensitive to the attochirp of the harmonic radiation, as illustrated in Fig.~{\ref{fig:1}}. The rescattered electron wave packet images different molecular landscapes depending on when it rescatters \cite{varju2005, niikura2005, smirnova2009, dudovich2006, azoury2017}, leading to a variation in the HHG light emission. High-frequency light (near the cutoff energy) is emitted later, meaning that an earlier delay is required to image any given position of the hole along the molecular backbone. Note that harmonic generation from any neutral molecules not undergoing CM would not have any delay dependence, and so would be canceled out by the normalization process. 

From the purple and green dashed lines in Fig.~{\ref{fig:1}}, we see that the HHG yield increases when the hole is located in the terminal bond (\emph{i.e.} when the electron density is on the bromine atom), and vice versa. This conclusion suggests that the scattering cross-section of the bromine atom is larger than the rest of the carbon chain, meaning that an increase in the overall density on the bromine atom (when the hole is \emph{not} on the halogen) results in a relative increase in the harmonic yield. This is a crucial result: because there is a spatially-resolvable feature in the harmonic spectrum -- here, a decrease in the harmonic yield when the hole is located on the halogen atom -- we are able to perform a time- and space-resolved measurement of the CM dynamics using frequency-matched strobo-spectroscopy (FMSS).

Though we are simulating and measuring particle-like CM dynamics \cite{folorunso2021, hamer2022} in this work, we expect that FMSS can be used to characterize a variety of ultrafast electron dynamics. The only requirement is that there is one or more features of the harmonic yield that can be traced back to specific parts of the molecule. As an example, in the usual way that CM is described, as a back-and-forth motion between two sites (\emph{e.g.} bromoacetylene), a measure of the amount of electron density on one of the sites fully describes the CM motion since any hole density not on the probed site must be on the other site. 

\begin{figure}[!t]
    \centering
    \includegraphics[width=0.65\linewidth]{Figures/figure4.png}
    \caption{(a) Normalized CM+FMSS spectrum, again for $\lambda_{L} = 1344\  \text{nm}$ and $I_{\circ} = 45\  \text{TW}/\text{cm}^{2}$, calculated from the dipole signal from the three $\pi_{xz}$ orbitals. Black dashed lines again taken from model calculations. (b) Amount of ionized charge, three optical cycles after the center of the laser pulse, as a function of delay.}
    \label{fig:4}
\end{figure}

While the CM orbital used in Fig.~{\ref{fig:3}} gives us the clearest picture of the CM dynamics (see again Fig.~{\ref{fig:2}}(b)), it does not correspond to a physical observable. Consider the electronic structure of $\text{BrC}_{4}\text{H}$: in addition to some lower-lying $\sigma$-type orbitals that do not contribute to the CM or the HHG, there are six $\pi$-type orbitals that span the length of the molecular backbone. Three of these $\pi$ orbitals lie in the $xz$-plane (where the molecular backbone is along the $z$-axis, and the laser is along the $x$-axis), while the other three lie in the $yz$-plane. By pulling one of the two electrons from one of the $\pi$ orbitals in the $xz$-plane (the CM orbital), we induce particle-like CM in $\text{BrC}_{4}\text{H}$; however, there are an additional 4 electrons in the $\pi_{xz}$ system that strongly contribute to both the CM and the HHG. Therefore, we look at the combined dipole signal from the three $\pi_{xz}$ orbitals. We have checked that these results are consistent with using the total dipole acceleration rather than the $\pi_{xz}$-orbitals-resolved dipole moment. 

Thus far, we have been looking at the relative increase in the delay-dependent HHG yield that occurs when the hole is $\emph{not}$ on the bromine atom. This method works well for the CM-orbital-resolved FMSS spectrum of Fig.~{\ref{fig:3}}(a); switching to the $\pi_{xz}$-orbitals-resolved FMSS spectrum, however, we instead look for an \emph{absence} of harmonic yield corresponding to the hole being on the bromine atom. Thus, in Fig.~{\ref{fig:4}}(a), we plot the inverse of the $\pi_{xz}$-orbitals-resolved harmonic yield, $1/S_{\pi}(\omega)$. Again, we see a delay- and harmonic-frequency-dependent variation in the (inverse) harmonic yield due to the CM dynamics. The black dashed lines, again taken from our model calculations in Fig.~{\ref{fig:3}}(b), have been shifted by $0.25\omega_{L}$ since we are looking for an absence, rather than the presence, of harmonic signal.

We have shown that the HHG yield tracks the hole density on the bromine atom. To further illustrate this, we algebraically remove the effect of the attochirp in the CM+FMSS spectra of Figs.~{\ref{fig:3}}(a) and \ref{fig:4}(a) in order to obtain an absolute-time-dependent measure of how much hole density is on the bromine end of the molecule. This analysis is performed in Figure~{\ref{fig:5}}. The blue curve depicts the amount of hole density centered around the bromine atom, taken from the field-free CM dynamics depicted in Fig.~{\ref{fig:2}}(b). We compare this hole density to the recombination-time-dependent harmonic yield, integrated over harmonic frequencies above 20 eV, for the CM-orbital-resolved data in Fig.~{\ref{fig:3}}(a) (solid red curve) and the $\pi_{xz}$-orbitals-resolved data in Fig.~{\ref{fig:4}}(a) (dashed red curve). From the semiclassical model of HHG \cite{corkum1993, schafer1993, lewenstein1994, chang2011}, we know exactly when each harmonic is emitted as a function of absolute time (for every delay $\tau$). From our TDDFT simulations, we also know the exact location of the electron hole as a function of absolute time. Thus, we can unambiguously map the variation in the harmonic signal to the amount of electron density on the halogen atom. In Fig.~{\ref{fig:5}}, a value near the top of the figure means the hole density is \emph{not} localized on the bromine atom (is localized on the terminal bond), and therefore results in a larger HHG yield. Despite the different methods used to obtain the red and blue curves in Fig.~{\ref{fig:5}}, they match each other very well. Note that the higher-frequency oscillations in the red curves (particularly, the dashed red curve) can be explained by the additional atomic-center-localized oscillations in the hole density in Fig.~{\ref{fig:2}}(c).

We finally ask the question of whether the CM-induced modulation of the harmonic yield is due to the ionization step or the rescattering step. To do so, in Figure~{\ref{fig:4}}(b) we plot the amount of charge ionized from the simulation box, one-half laser cycle after the end of the laser pulse, as a function of the delay $\tau$. There is a small amount of leakage -- ionized charge leaving the simulation box even in the absence of the laser field, approximately 2\% of an electron per laser cycle -- as evidenced by the overall slope in Fig.~{\ref{fig:4}}(b), which can be attributed to the absorbing boundaries in the direction perpendicular to both the CM and the laser. We have checked that the leakage does not effect the results shown here. On top of this overall linear slope, we see a clear a half-laser-cycle-periodic modulation in the ionization signal due to the CM dynamics. Different relative phases between the CM and the peaks of the laser field cause different amounts of charge to be ionized as a function of $\tau$. However, after correcting for the leakage, the amplitude of the oscillation in the ionization signal is quite small (roughly 1\%) compared to the variation in the harmonic signal, for which the yield is roughly three times larger when the hole is not on the bromine atom (as opposed to five times larger, for the CM-orbital-resolved case). Thus, we conclude that the modulation of the HHG signal from the CM dynamics occurs mainly as a result of the recombination step, not the ionization step.

\begin{figure}[!t]
    \centering
    \includegraphics[width=0.65\linewidth]{Figures/figure5.png}
    \caption{Comparison between time-dependent, field-free hole density on the bromine atom (blue) and recombination-time-dependent HHG yield, with the attochirp removed and integrated over harmonic frequencies above $20\ \text{eV}$ (red). The solid red line is taken from the CM-orbital-resolved spectra in Fig.~{\ref{fig:3}(a)}, and the dashed line corresponds to the $\pi_{xz}$-orbitals-resolved data in Fig.~{\ref{fig:4}(a)}.}
    \label{fig:5}
\end{figure}

\section{Conclusions}

\noindent In summary, we have shown that frequency-matched high-harmonic strobo-spectroscopy of charge migration (CM+FMSS) in $\text{BrC}_{4}\text{H}$ causes a coherent modulation of the HHG signal that precisely tracks the amount of electron density on the bromine atom, which tells us the phase of the CM motion. By exploiting a site-specific feature of the HHG spectrum, we achieve a time- and space-resolved measurement of the CM by performing a sub-cycle-resolved delay scan. FMSS takes advantage of the intrinsic attosecond time resolution of the HHG process (the attochirp), in which different harmonics are emitted at different times and thus probe different locations of the electron hole. These claims are supported by a similar result from an SFA-inspired model calculation. We can also make a direct comparison between the recombination-time-dependent, harmonic-frequency-integrated HHG yield and the hole density on the halogen.

Beyond the $\text{BrC}_{4}\text{H}$ molecule used here, we note that similar particle-like CM dynamics have been predicted in other classes of molecules \cite{folorunso2021, folorunso2023}. Thus, given the generalizable nature of our approach, we expect that CM+FMSS analyses can be applied broadly to other classes of molecules, such as functionalized benzenes or even bio-molecules and beyond.

%Using these results, we present a method to measure CM dynamics of an unknown molecular system with high-harmonic spectroscopy: first, impulsively align the sample and initiate CM dynamics using strong-field ionization. Then, irradiate the molecule using a wavelength-tunable, many-cycle laser pulse to drive HHG. Assuming that the applied laser frequency and the fundamental CM frequency are incommensurate, sidebands will appear in the harmonic spectrum due to the molecules that are undergoing CM \cite{hamer2022}. Since the location of these sidebands changes when scanning the wavelength of the applied laser, one can make a precise measurement of the CM frequency \cite{hamer2022}. Now that we know the CM frequency, we can irradiate the sample with a frequency-matched ($\omega_{L} = \omega_{\text{CM}}/2$) few-cycle laser pulse, and scan the delay between the initiation of the CM dynamics and the center of the laser pulse to obtain the phase of the CM dynamics. Note that attempting a delay scan using the frequency-unmatched many-cycle pulse above would result in a delay-independent signal since the electron hole would be at different positions along the molecular backbone for each half-cycle of the laser field, and a many-cycle pulse would average over all of those positions. This suggests that a sufficiently short laser pulse would allow for some inaccuracy in the measurement of the CM frequency since this delay-averaging would be much less severe. From these two measurements alone, we have fully characterized the CM motion in the molecule using high-harmonic spectroscopy.

\section{Supporting Information}

The TDDFT simulation data and Python scripts we use to produce the figures are available at [authors will provide link to public repository for production].

\section{Acknowledgments}

We thank L. F. DiMauro and R. R. Jones for helpful discussions on this work.
%
This work was supported by the U.S.\ Department of Energy, Office of Science, Office of Basic Energy Sciences, under Award No.~DE-SC0012462.
%
Portions of this research were conducted with high performance computational resources provided by Louisiana State University (\url{http://www.hpc.lsu.edu}) and the Louisiana Optical Network Infrastructure (\url{http://www.loni.org}).

\singlespacing
\providecommand{\latin}[1]{#1}
\makeatletter
\providecommand{\doi}
  {\begingroup\let\do\@makeother\dospecials
  \catcode`\{=1 \catcode`\}=2 \doi@aux}
\providecommand{\doi@aux}[1]{\endgroup\texttt{#1}}
\makeatother
\providecommand*\mcitethebibliography{\thebibliography}
\csname @ifundefined\endcsname{endmcitethebibliography}
  {\let\endmcitethebibliography\endthebibliography}{}
\begin{mcitethebibliography}{61}
\providecommand*\natexlab[1]{#1}
\providecommand*\mciteSetBstSublistMode[1]{}
\providecommand*\mciteSetBstMaxWidthForm[2]{}
\providecommand*\mciteBstWouldAddEndPuncttrue
  {\def\EndOfBibitem{\unskip.}}
\providecommand*\mciteBstWouldAddEndPunctfalse
  {\let\EndOfBibitem\relax}
\providecommand*\mciteSetBstMidEndSepPunct[3]{}
\providecommand*\mciteSetBstSublistLabelBeginEnd[3]{}
\providecommand*\EndOfBibitem{}
\mciteSetBstSublistMode{f}
\mciteSetBstMaxWidthForm{subitem}{(\alph{mcitesubitemcount})}
\mciteSetBstSublistLabelBeginEnd
  {\mcitemaxwidthsubitemform\space}
  {\relax}
  {\relax}

\bibitem[Calegari \latin{et~al.}(2016)Calegari, Trabattoni, Palacios, Ayuso,
  Castrovilli, Greenwood, Decleva, Mart{\'{\i}}n, and Nisoli]{calegari2016}
Calegari,~F.; Trabattoni,~A.; Palacios,~A.; Ayuso,~D.; Castrovilli,~M.~C.;
  Greenwood,~J.~B.; Decleva,~P.; Mart{\'{\i}}n,~F.; Nisoli,~M. Charge migration
  induced by attosecond pulses in bio-relevant molecules. \emph{Journal of
  Physics B: Atomic, Molecular and Optical Physics} \textbf{2016}, \emph{49},
  142001\relax
\mciteBstWouldAddEndPuncttrue
\mciteSetBstMidEndSepPunct{\mcitedefaultmidpunct}
{\mcitedefaultendpunct}{\mcitedefaultseppunct}\relax
\EndOfBibitem
\bibitem[Wörner \latin{et~al.}(2017)Wörner, Arrell, Banerji, Cannizzo,
  Chergui, Das, Hamm, Keller, Kraus, Liberatore, Lopez-Tarifa, Lucchini,
  Meuwly, Milne, Moser, Rothlisberger, Smolentsev, Teuscher, van Bokhoven, and
  Wenger]{worner2017}
Wörner,~H.~J. \latin{et~al.}  Charge migration and charge transfer in
  molecular systems. \emph{Structural Dynamics} \textbf{2017}, \emph{4},
  061508\relax
\mciteBstWouldAddEndPuncttrue
\mciteSetBstMidEndSepPunct{\mcitedefaultmidpunct}
{\mcitedefaultendpunct}{\mcitedefaultseppunct}\relax
\EndOfBibitem
\bibitem[Nisoli \latin{et~al.}(2017)Nisoli, Decleva, Calegari, Palacios, and
  Mart\'{\i}n]{nisoli2017}
Nisoli,~M.; Decleva,~P.; Calegari,~F.; Palacios,~A.; Mart\'{\i}n,~F. Attosecond
  Electron Dynamics in Molecules. \emph{Chemical Reviews} \textbf{2017},
  \emph{117}, 10760--10825\relax
\mciteBstWouldAddEndPuncttrue
\mciteSetBstMidEndSepPunct{\mcitedefaultmidpunct}
{\mcitedefaultendpunct}{\mcitedefaultseppunct}\relax
\EndOfBibitem
\bibitem[Li \latin{et~al.}(2020)Li, Lu, Chew, Han, Li, Wu, Wang, Ghimire, and
  Chang]{li2020}
Li,~J.; Lu,~J.; Chew,~A.; Han,~S.; Li,~J.; Wu,~Y.; Wang,~H.; Ghimire,~S.;
  Chang,~Z. Attosecond science based on high harmonic generation from gases and
  solids. \emph{Nature Communications} \textbf{2020}, \emph{11}, 2748\relax
\mciteBstWouldAddEndPuncttrue
\mciteSetBstMidEndSepPunct{\mcitedefaultmidpunct}
{\mcitedefaultendpunct}{\mcitedefaultseppunct}\relax
\EndOfBibitem
\bibitem[Goulielmakis \latin{et~al.}(2010)Goulielmakis, Loh, Wirth, Santra,
  Rohringer, Yakovlev, Zherebtsov, Pfeifer, Azzeer, Kling, Leone, and
  Krausz]{goulielmakis2010}
Goulielmakis,~E.; Loh,~Z.-H.; Wirth,~A.; Santra,~R.; Rohringer,~N.;
  Yakovlev,~V.~S.; Zherebtsov,~S.; Pfeifer,~T.; Azzeer,~A.~M.; Kling,~M.~F.;
  Leone,~S.~R.; Krausz,~F. Real-time observation of valence electron motion.
  \emph{Nature} \textbf{2010}, \emph{466}, 739--743\relax
\mciteBstWouldAddEndPuncttrue
\mciteSetBstMidEndSepPunct{\mcitedefaultmidpunct}
{\mcitedefaultendpunct}{\mcitedefaultseppunct}\relax
\EndOfBibitem
\bibitem[Calegari \latin{et~al.}(2014)Calegari, Ayuso, Trabattoni, Belshaw,
  Camillis, Anumula, Frassetto, Poletto, Palacios, Decleva, Greenwood,
  Mart\'{\i}n, and Nisoli]{calegari2014}
Calegari,~F.; Ayuso,~D.; Trabattoni,~A.; Belshaw,~L.; Camillis,~S.~D.;
  Anumula,~S.; Frassetto,~F.; Poletto,~L.; Palacios,~A.; Decleva,~P.;
  Greenwood,~J.~B.; Mart\'{\i}n,~F.; Nisoli,~M. Ultrafast electron dynamics in
  phenylalanine initiated by attosecond pulses. \emph{Science} \textbf{2014},
  \emph{346}, 336--339\relax
\mciteBstWouldAddEndPuncttrue
\mciteSetBstMidEndSepPunct{\mcitedefaultmidpunct}
{\mcitedefaultendpunct}{\mcitedefaultseppunct}\relax
\EndOfBibitem
\bibitem[Kraus \latin{et~al.}(2015)Kraus, Mignolet, Baykusheva, Rupenyan,
  Horn\'{y}, Penka, Grassi, Tolstikhin, Schneider, Jensen, Madsen, Bandrauk,
  Remacle, and Wörner]{kraus2015}
Kraus,~P.~M.; Mignolet,~B.; Baykusheva,~D.; Rupenyan,~A.; Horn\'{y},~L.;
  Penka,~E.~F.; Grassi,~G.; Tolstikhin,~O.~I.; Schneider,~J.; Jensen,~F.;
  Madsen,~L.~B.; Bandrauk,~A.~D.; Remacle,~F.; Wörner,~H.~J. Measurement and
  laser control of attosecond charge migration in ionized iodoacetylene.
  \emph{Science} \textbf{2015}, \emph{350}, 790--795\relax
\mciteBstWouldAddEndPuncttrue
\mciteSetBstMidEndSepPunct{\mcitedefaultmidpunct}
{\mcitedefaultendpunct}{\mcitedefaultseppunct}\relax
\EndOfBibitem
\bibitem[Remacle \latin{et~al.}(1998)Remacle, Levine, and Ratner]{remacle1998}
Remacle,~F.; Levine,~R.; Ratner,~M. Charge directed reactivity: a simple
  electronic model, exhibiting site selectivity, for the dissociation of ions.
  \emph{Chemical Physics Letters} \textbf{1998}, \emph{285}, 25--33\relax
\mciteBstWouldAddEndPuncttrue
\mciteSetBstMidEndSepPunct{\mcitedefaultmidpunct}
{\mcitedefaultendpunct}{\mcitedefaultseppunct}\relax
\EndOfBibitem
\bibitem[Golubev and Kuleff(2015)Golubev, and Kuleff]{golubev2015}
Golubev,~N.~V.; Kuleff,~A.~I. Control of charge migration in molecules by
  ultrashort laser pulses. \emph{Phys. Rev. A} \textbf{2015}, \emph{91},
  051401\relax
\mciteBstWouldAddEndPuncttrue
\mciteSetBstMidEndSepPunct{\mcitedefaultmidpunct}
{\mcitedefaultendpunct}{\mcitedefaultseppunct}\relax
\EndOfBibitem
\bibitem[Mauger \latin{et~al.}(2022)Mauger, Folorunso, Hamer, Chandre, Gaarde,
  Lopata, and Schafer]{mauger2022}
Mauger,~F.; Folorunso,~A.~S.; Hamer,~K.~A.; Chandre,~C.; Gaarde,~M.~B.;
  Lopata,~K.; Schafer,~K.~J. Charge migration and attosecond solitons in
  conjugated organic molecules. \emph{Phys. Rev. Research} \textbf{2022},
  \emph{4}, 013073\relax
\mciteBstWouldAddEndPuncttrue
\mciteSetBstMidEndSepPunct{\mcitedefaultmidpunct}
{\mcitedefaultendpunct}{\mcitedefaultseppunct}\relax
\EndOfBibitem
\bibitem[Cederbaum and Zobeley(1999)Cederbaum, and Zobeley]{cederbaum1999}
Cederbaum,~L.; Zobeley,~J. Ultrafast charge migration by electron correlation.
  \emph{Chemical Physics Letters} \textbf{1999}, \emph{307}, 205--210\relax
\mciteBstWouldAddEndPuncttrue
\mciteSetBstMidEndSepPunct{\mcitedefaultmidpunct}
{\mcitedefaultendpunct}{\mcitedefaultseppunct}\relax
\EndOfBibitem
\bibitem[Breidbach and Cederbaum(2003)Breidbach, and Cederbaum]{breidbach2003}
Breidbach,~J.; Cederbaum,~L.~S. Migration of holes: Formalism, mechanisms, and
  illustrative applications. \emph{The Journal of Chemical Physics}
  \textbf{2003}, \emph{118}, 3983--3996\relax
\mciteBstWouldAddEndPuncttrue
\mciteSetBstMidEndSepPunct{\mcitedefaultmidpunct}
{\mcitedefaultendpunct}{\mcitedefaultseppunct}\relax
\EndOfBibitem
\bibitem[Lünnemann \latin{et~al.}(2008)Lünnemann, Kuleff, and
  Cederbaum]{lunnemann2008}
Lünnemann,~S.; Kuleff,~A.~I.; Cederbaum,~L.~S. Ultrafast charge migration in
  2-phenylethyl-N,N-dimethylamine. \emph{Chemical Physics Letters}
  \textbf{2008}, \emph{450}, 232--235\relax
\mciteBstWouldAddEndPuncttrue
\mciteSetBstMidEndSepPunct{\mcitedefaultmidpunct}
{\mcitedefaultendpunct}{\mcitedefaultseppunct}\relax
\EndOfBibitem
\bibitem[Kuleff \latin{et~al.}(2016)Kuleff, Kryzhevoi, Pernpointner, and
  Cederbaum]{kuleff2016}
Kuleff,~A.~I.; Kryzhevoi,~N.~V.; Pernpointner,~M.; Cederbaum,~L.~S. Core
  Ionization Initiates Subfemtosecond Charge Migration in the Valence Shell of
  Molecules. \emph{Phys. Rev. Lett.} \textbf{2016}, \emph{117}, 093002\relax
\mciteBstWouldAddEndPuncttrue
\mciteSetBstMidEndSepPunct{\mcitedefaultmidpunct}
{\mcitedefaultendpunct}{\mcitedefaultseppunct}\relax
\EndOfBibitem
\bibitem[Bruner \latin{et~al.}(2017)Bruner, Hernandez, Mauger, Abanador,
  LaMaster, Gaarde, Schafer, and Lopata]{bruner2017}
Bruner,~A.; Hernandez,~S.; Mauger,~F.; Abanador,~P.~M.; LaMaster,~D.~J.;
  Gaarde,~M.~B.; Schafer,~K.~J.; Lopata,~K. Attosecond Charge Migration with
  TDDFT: Accurate Dynamics from a Well-Defined Initial State. \emph{The Journal
  of Physical Chemistry Letters} \textbf{2017}, \emph{8}, 3991--3996\relax
\mciteBstWouldAddEndPuncttrue
\mciteSetBstMidEndSepPunct{\mcitedefaultmidpunct}
{\mcitedefaultendpunct}{\mcitedefaultseppunct}\relax
\EndOfBibitem
\bibitem[Folorunso \latin{et~al.}(2021)Folorunso, Bruner, Mauger, Hamer,
  Hernandez, Jones, DiMauro, Gaarde, Schafer, and Lopata]{folorunso2021}
Folorunso,~A.~S.; Bruner,~A.; Mauger,~F.; Hamer,~K.~A.; Hernandez,~S.;
  Jones,~R.~R.; DiMauro,~L.~F.; Gaarde,~M.~B.; Schafer,~K.~J.; Lopata,~K.
  Molecular Modes of Attosecond Charge Migration. \emph{Phys. Rev. Lett.}
  \textbf{2021}, \emph{126}, 133002\relax
\mciteBstWouldAddEndPuncttrue
\mciteSetBstMidEndSepPunct{\mcitedefaultmidpunct}
{\mcitedefaultendpunct}{\mcitedefaultseppunct}\relax
\EndOfBibitem
\bibitem[M{\aa}nsson \latin{et~al.}(2021)M{\aa}nsson, Latini, Covito, Wanie,
  Galli, Perfetto, Stefanucci, H{\"u}bener, De~Giovannini, Castrovilli,
  Trabattoni, Frassetto, Poletto, Greenwood, L{\'e}gar{\'e}, Nisoli, Rubio, and
  Calegari]{mansson2021}
M{\aa}nsson,~E.~P. \latin{et~al.}  Real-time observation of a
  correlation-driven sub 3 fs charge migration in ionised adenine.
  \emph{Communications Chemistry} \textbf{2021}, \emph{4}, 73\relax
\mciteBstWouldAddEndPuncttrue
\mciteSetBstMidEndSepPunct{\mcitedefaultmidpunct}
{\mcitedefaultendpunct}{\mcitedefaultseppunct}\relax
\EndOfBibitem
\bibitem[Barillot \latin{et~al.}(2021)Barillot, Alexander, Cooper, Driver,
  Garratt, Li, Al~Haddad, Sanchez-Gonzalez, Ag\aa{}ker, Arrell, Bearpark,
  Berrah, Bostedt, Bozek, Brahms, Bucksbaum, Clark, Doumy, Feifel, Frasinski,
  Jarosch, Johnson, Kjellsson, Koloren\c{c}, Kumagai, Larsen, Matia-Hernando,
  Robb, Rubensson, Ruberti, Sathe, Squibb, Tan, Tisch, Vacher, Walke, Wolf,
  Wood, Zhaunerchyk, Walter, Osipov, Marinelli, Maxwell, Coffee, Lutman,
  Averbukh, Ueda, Cryan, and Marangos]{barillot2022}
Barillot,~T. \latin{et~al.}  Correlation-Driven Transient Hole Dynamics
  Resolved in Space and Time in the Isopropanol Molecule. \emph{Phys. Rev. X}
  \textbf{2021}, \emph{11}, 031048\relax
\mciteBstWouldAddEndPuncttrue
\mciteSetBstMidEndSepPunct{\mcitedefaultmidpunct}
{\mcitedefaultendpunct}{\mcitedefaultseppunct}\relax
\EndOfBibitem
\bibitem[Li \latin{et~al.}(2022)Li, Driver, Rosenberger, Champenois, Duris,
  Al-Haddad, Averbukh, Barnard, Berrah, Bostedt, Bucksbaum, Coffee, DiMauro,
  Fang, Garratt, Gatton, Guo, Hartmann, Haxton, Helml, Huang, LaForge, Kamalov,
  Knurr, Lin, Lutman, MacArthur, Marangos, Nantel, Natan, Obaid, O’Neal,
  Shivaram, Schori, Walter, Wang, Wolf, Zhang, Kling, Marinelli, and
  Cryan]{li2022}
Li,~S. \latin{et~al.}  Attosecond coherent electron motion in Auger-Meitner
  decay. \emph{Science} \textbf{2022}, \emph{375}, 285--290\relax
\mciteBstWouldAddEndPuncttrue
\mciteSetBstMidEndSepPunct{\mcitedefaultmidpunct}
{\mcitedefaultendpunct}{\mcitedefaultseppunct}\relax
\EndOfBibitem
\bibitem[Schlegel \latin{et~al.}(2022)Schlegel, Hoerner, and Li]{schlegel2022}
Schlegel,~H.~B.; Hoerner,~P.; Li,~W. Ionization of HCCI Neutral and Cations by
  Strong Laser Fields Simulated With Time Dependent Configuration Interaction.
  \emph{Frontiers in Chemistry} \textbf{2022}, \emph{10}\relax
\mciteBstWouldAddEndPuncttrue
\mciteSetBstMidEndSepPunct{\mcitedefaultmidpunct}
{\mcitedefaultendpunct}{\mcitedefaultseppunct}\relax
\EndOfBibitem
\bibitem[Yong \latin{et~al.}(2022)Yong, Sun, Gu, and Mukamel]{yong2022}
Yong,~H.; Sun,~S.; Gu,~B.; Mukamel,~S. Attosecond Charge Migration in Molecules
  Imaged by Combined X-ray and Electron Diffraction. \emph{Journal of the
  American Chemical Society} \textbf{2022}, \emph{144}, 20710--20716\relax
\mciteBstWouldAddEndPuncttrue
\mciteSetBstMidEndSepPunct{\mcitedefaultmidpunct}
{\mcitedefaultendpunct}{\mcitedefaultseppunct}\relax
\EndOfBibitem
\bibitem[Zhao \latin{et~al.}(2022)Zhao, Liu, Wang, Yuan, and Zhao]{zhao2022}
Zhao,~J.; Liu,~J.; Wang,~X.; Yuan,~J.; Zhao,~Z. Real-Time Observation of
  Electron-Hole Coherence Induced by Strong-Field Ionization. \emph{Chinese
  Physics Letters} \textbf{2022}, \emph{39}, 123201\relax
\mciteBstWouldAddEndPuncttrue
\mciteSetBstMidEndSepPunct{\mcitedefaultmidpunct}
{\mcitedefaultendpunct}{\mcitedefaultseppunct}\relax
\EndOfBibitem
\bibitem[Folorunso \latin{et~al.}(2023)Folorunso, Mauger, Hamer, Jayasinghe,
  Wahyutama, Ragains, Jones, DiMauro, Gaarde, Schafer, and
  Lopata]{folorunso2023}
Folorunso,~A.~S.; Mauger,~F.; Hamer,~K.~A.; Jayasinghe,~D.~D.;
  Wahyutama,~I.~S.; Ragains,~J.~R.; Jones,~R.~R.; DiMauro,~L.~F.;
  Gaarde,~M.~B.; Schafer,~K.~J.; Lopata,~K. Attochemistry Regulation of Charge
  Migration. \emph{The Journal of Physical Chemistry A} \textbf{2023},
  \emph{127}, 1894--1900\relax
\mciteBstWouldAddEndPuncttrue
\mciteSetBstMidEndSepPunct{\mcitedefaultmidpunct}
{\mcitedefaultendpunct}{\mcitedefaultseppunct}\relax
\EndOfBibitem
\bibitem[Yu \latin{et~al.}(2023)Yu, Liang, Geng, and Peng]{yu2023}
Yu,~W.-D.; Liang,~H.; Geng,~L.; Peng,~L.-Y. Dynamical analysis of attosecond
  molecular modes. \emph{Phys. Rev. A} \textbf{2023}, \emph{107}, 013101\relax
\mciteBstWouldAddEndPuncttrue
\mciteSetBstMidEndSepPunct{\mcitedefaultmidpunct}
{\mcitedefaultendpunct}{\mcitedefaultseppunct}\relax
\EndOfBibitem
\bibitem[He \latin{et~al.}(2022)He, Sun, Lan, He, Wang, Wang, Zhu, Li, Cao, Lu,
  and Lin]{he2022}
He,~L.; Sun,~S.; Lan,~P.; He,~Y.; Wang,~B.; Wang,~P.; Zhu,~X.; Li,~L.; Cao,~W.;
  Lu,~P.; Lin,~C.~D. Filming movies of attosecond charge migration in single
  molecules with high harmonic spectroscopy. \emph{Nature Communications}
  \textbf{2022}, \emph{13}, 4595\relax
\mciteBstWouldAddEndPuncttrue
\mciteSetBstMidEndSepPunct{\mcitedefaultmidpunct}
{\mcitedefaultendpunct}{\mcitedefaultseppunct}\relax
\EndOfBibitem
\bibitem[Hamer \latin{et~al.}(2022)Hamer, Mauger, Folorunso, Lopata, Jones,
  DiMauro, Schafer, and Gaarde]{hamer2022}
Hamer,~K.~A.; Mauger,~F.; Folorunso,~A.~S.; Lopata,~K.; Jones,~R.~R.;
  DiMauro,~L.~F.; Schafer,~K.~J.; Gaarde,~M.~B. Characterizing particle-like
  charge-migration dynamics with high-order harmonic sideband spectroscopy.
  \emph{Phys. Rev. A} \textbf{2022}, \emph{106}, 013103\relax
\mciteBstWouldAddEndPuncttrue
\mciteSetBstMidEndSepPunct{\mcitedefaultmidpunct}
{\mcitedefaultendpunct}{\mcitedefaultseppunct}\relax
\EndOfBibitem
\bibitem[Matselyukh \latin{et~al.}(2022)Matselyukh, Despr{\'e}, Golubev,
  Kuleff, and W{\"o}rner]{matselyukh2022}
Matselyukh,~D.~T.; Despr{\'e},~V.; Golubev,~N.~V.; Kuleff,~A.~I.;
  W{\"o}rner,~H.~J. Decoherence and revival in attosecond charge migration
  driven by non-adiabatic dynamics. \emph{Nature Physics} \textbf{2022},
  \emph{18}, 1206--1213\relax
\mciteBstWouldAddEndPuncttrue
\mciteSetBstMidEndSepPunct{\mcitedefaultmidpunct}
{\mcitedefaultendpunct}{\mcitedefaultseppunct}\relax
\EndOfBibitem
\bibitem[Kobayashi \latin{et~al.}(2022)Kobayashi, Neumark, and
  Leone]{kobayashi2022}
Kobayashi,~Y.; Neumark,~D.~M.; Leone,~S.~R. Theoretical analysis of the role of
  complex transition dipole phase in XUV transient-absorption probing of charge
  migration. \emph{Opt. Express} \textbf{2022}, \emph{30}, 5673--5682\relax
\mciteBstWouldAddEndPuncttrue
\mciteSetBstMidEndSepPunct{\mcitedefaultmidpunct}
{\mcitedefaultendpunct}{\mcitedefaultseppunct}\relax
\EndOfBibitem
\bibitem[Bressler and Chergui(2010)Bressler, and Chergui]{bressler2010}
Bressler,~C.; Chergui,~M. Molecular Structural Dynamics Probed by Ultrafast
  X-Ray Absorption Spectroscopy. \emph{Annual Review of Physical Chemistry}
  \textbf{2010}, \emph{61}, 263--282\relax
\mciteBstWouldAddEndPuncttrue
\mciteSetBstMidEndSepPunct{\mcitedefaultmidpunct}
{\mcitedefaultendpunct}{\mcitedefaultseppunct}\relax
\EndOfBibitem
\bibitem[Bandrauk \latin{et~al.}(2004)Bandrauk, Chelkowski, and
  Nguyen]{bandrauk2004}
Bandrauk,~A.~D.; Chelkowski,~S.; Nguyen,~H.~S. Attosecond localization of
  electrons in molecules. \emph{International Journal of Quantum Chemistry}
  \textbf{2004}, \emph{100}, 834--844\relax
\mciteBstWouldAddEndPuncttrue
\mciteSetBstMidEndSepPunct{\mcitedefaultmidpunct}
{\mcitedefaultendpunct}{\mcitedefaultseppunct}\relax
\EndOfBibitem
\bibitem[Stolow \latin{et~al.}(2004)Stolow, Bragg, and Neumark]{stolow2004}
Stolow,~A.; Bragg,~A.~E.; Neumark,~D.~M. Femtosecond Time-Resolved
  Photoelectron Spectroscopy. \emph{Chemical Reviews} \textbf{2004},
  \emph{104}, 1719--1758\relax
\mciteBstWouldAddEndPuncttrue
\mciteSetBstMidEndSepPunct{\mcitedefaultmidpunct}
{\mcitedefaultendpunct}{\mcitedefaultseppunct}\relax
\EndOfBibitem
\bibitem[Ayuso \latin{et~al.}(2017)Ayuso, Palacios, Decleva, and
  Mart\'{\i}n]{ayuso2017}
Ayuso,~D.; Palacios,~A.; Decleva,~P.; Mart\'{\i}n,~F. Ultrafast charge dynamics
  in glycine induced by attosecond pulses. \emph{Phys. Chem. Chem. Phys.}
  \textbf{2017}, \emph{19}, 19767--19776\relax
\mciteBstWouldAddEndPuncttrue
\mciteSetBstMidEndSepPunct{\mcitedefaultmidpunct}
{\mcitedefaultendpunct}{\mcitedefaultseppunct}\relax
\EndOfBibitem
\bibitem[Lara-Astiaso \latin{et~al.}(2018)Lara-Astiaso, Galli, Trabattoni,
  Palacios, Ayuso, Frassetto, Poletto, De~Camillis, Greenwood, Decleva,
  Tavernelli, Calegari, Nisoli, and Mart\'{\i}n]{lara-astiaso2018}
Lara-Astiaso,~M.; Galli,~M.; Trabattoni,~A.; Palacios,~A.; Ayuso,~D.;
  Frassetto,~F.; Poletto,~L.; De~Camillis,~S.; Greenwood,~J.; Decleva,~P.;
  Tavernelli,~I.; Calegari,~F.; Nisoli,~M.; Mart\'{\i}n,~F. Attosecond
  Pump–Probe Spectroscopy of Charge Dynamics in Tryptophan. \emph{The Journal
  of Physical Chemistry Letters} \textbf{2018}, \emph{9}, 4570--4577\relax
\mciteBstWouldAddEndPuncttrue
\mciteSetBstMidEndSepPunct{\mcitedefaultmidpunct}
{\mcitedefaultendpunct}{\mcitedefaultseppunct}\relax
\EndOfBibitem
\bibitem[Bredtmann \latin{et~al.}(2012)Bredtmann, Chelkowski, and
  Bandrauk]{bredtmann2012}
Bredtmann,~T.; Chelkowski,~S.; Bandrauk,~A.~D. Effect of Nuclear Motion on
  Molecular High Order Harmonic Pump Probe Spectroscopy. \emph{The Journal of
  Physical Chemistry A} \textbf{2012}, \emph{116}, 11398--11405\relax
\mciteBstWouldAddEndPuncttrue
\mciteSetBstMidEndSepPunct{\mcitedefaultmidpunct}
{\mcitedefaultendpunct}{\mcitedefaultseppunct}\relax
\EndOfBibitem
\bibitem[Kraus and Wörner(2013)Kraus, and Wörner]{worner2013}
Kraus,~P.~M.; Wörner,~H.~J. Time-resolved high-harmonic spectroscopy of
  valence electron dynamics. \emph{Chemical Physics} \textbf{2013}, \emph{414},
  32--44\relax
\mciteBstWouldAddEndPuncttrue
\mciteSetBstMidEndSepPunct{\mcitedefaultmidpunct}
{\mcitedefaultendpunct}{\mcitedefaultseppunct}\relax
\EndOfBibitem
\bibitem[Peng \latin{et~al.}(2019)Peng, Marceau, and Villeneuve]{peng2019}
Peng,~P.; Marceau,~C.; Villeneuve,~D.~M. Attosecond imaging of molecules using
  high harmonic spectroscopy. \emph{Nature Reviews Physics} \textbf{2019},
  \emph{1}, 144--155\relax
\mciteBstWouldAddEndPuncttrue
\mciteSetBstMidEndSepPunct{\mcitedefaultmidpunct}
{\mcitedefaultendpunct}{\mcitedefaultseppunct}\relax
\EndOfBibitem
\bibitem[Dudovich \latin{et~al.}(2006)Dudovich, Smirnova, Levesque, Mairesse,
  Ivanov, Villeneuve, and Corkum]{dudovich2006}
Dudovich,~N.; Smirnova,~O.; Levesque,~J.; Mairesse,~Y.; Ivanov,~M.~Y.;
  Villeneuve,~D.~M.; Corkum,~P.~B. Measuring and controlling the birth of
  attosecond XUV pulses. \emph{Nature Physics} \textbf{2006}, \emph{2},
  781--786\relax
\mciteBstWouldAddEndPuncttrue
\mciteSetBstMidEndSepPunct{\mcitedefaultmidpunct}
{\mcitedefaultendpunct}{\mcitedefaultseppunct}\relax
\EndOfBibitem
\bibitem[Smirnova \latin{et~al.}(2009)Smirnova, Patchkovskii, Mairesse,
  Dudovich, and Ivanov]{smirnova2009}
Smirnova,~O.; Patchkovskii,~S.; Mairesse,~Y.; Dudovich,~N.; Ivanov,~M.~Y.
  Strong-field control and spectroscopy of attosecond electron-hole dynamics in
  molecules. \emph{Proceedings of the National Academy of Sciences}
  \textbf{2009}, \emph{106}, 16556--16561\relax
\mciteBstWouldAddEndPuncttrue
\mciteSetBstMidEndSepPunct{\mcitedefaultmidpunct}
{\mcitedefaultendpunct}{\mcitedefaultseppunct}\relax
\EndOfBibitem
\bibitem[Azoury \latin{et~al.}(2017)Azoury, Kr{\"u}ger, Orenstein, Larsson,
  Bauch, Bruner, and Dudovich]{azoury2017}
Azoury,~D.; Kr{\"u}ger,~M.; Orenstein,~G.; Larsson,~H.~R.; Bauch,~S.;
  Bruner,~B.~D.; Dudovich,~N. Self-probing spectroscopy of XUV photo-ionization
  dynamics in atoms subjected to a strong-field environment. \emph{Nature
  Communications} \textbf{2017}, \emph{8}, 1453\relax
\mciteBstWouldAddEndPuncttrue
\mciteSetBstMidEndSepPunct{\mcitedefaultmidpunct}
{\mcitedefaultendpunct}{\mcitedefaultseppunct}\relax
\EndOfBibitem
\bibitem[Kohn and Sham(1965)Kohn, and Sham]{kohn1965}
Kohn,~W.; Sham,~L.~J. Self-Consistent Equations Including Exchange and
  Correlation Effects. \emph{Phys. Rev.} \textbf{1965}, \emph{140},
  A1133--A1138\relax
\mciteBstWouldAddEndPuncttrue
\mciteSetBstMidEndSepPunct{\mcitedefaultmidpunct}
{\mcitedefaultendpunct}{\mcitedefaultseppunct}\relax
\EndOfBibitem
\bibitem[Runge and Gross(1984)Runge, and Gross]{runge1984}
Runge,~E.; Gross,~E. K.~U. Density-Functional Theory for Time-Dependent
  Systems. \emph{Phys. Rev. Lett.} \textbf{1984}, \emph{52}, 997--1000\relax
\mciteBstWouldAddEndPuncttrue
\mciteSetBstMidEndSepPunct{\mcitedefaultmidpunct}
{\mcitedefaultendpunct}{\mcitedefaultseppunct}\relax
\EndOfBibitem
\bibitem[Corkum(1993)]{corkum1993}
Corkum,~P.~B. Plasma perspective on strong field multiphoton ionization.
  \emph{Phys. Rev. Lett.} \textbf{1993}, \emph{71}, 1994--1997\relax
\mciteBstWouldAddEndPuncttrue
\mciteSetBstMidEndSepPunct{\mcitedefaultmidpunct}
{\mcitedefaultendpunct}{\mcitedefaultseppunct}\relax
\EndOfBibitem
\bibitem[Schafer \latin{et~al.}(1993)Schafer, Yang, DiMauro, and
  Kulander]{schafer1993}
Schafer,~K.~J.; Yang,~B.; DiMauro,~L.~F.; Kulander,~K.~C. Above threshold
  ionization beyond the high harmonic cutoff. \emph{Phys. Rev. Lett.}
  \textbf{1993}, \emph{70}, 1599--1602\relax
\mciteBstWouldAddEndPuncttrue
\mciteSetBstMidEndSepPunct{\mcitedefaultmidpunct}
{\mcitedefaultendpunct}{\mcitedefaultseppunct}\relax
\EndOfBibitem
\bibitem[Lewenstein \latin{et~al.}(1994)Lewenstein, Balcou, Ivanov, L'Huillier,
  and Corkum]{lewenstein1994}
Lewenstein,~M.; Balcou,~P.; Ivanov,~M.~Y.; L'Huillier,~A.; Corkum,~P.~B. Theory
  of high-harmonic generation by low-frequency laser fields. \emph{Phys. Rev.
  A} \textbf{1994}, \emph{49}, 2117--2132\relax
\mciteBstWouldAddEndPuncttrue
\mciteSetBstMidEndSepPunct{\mcitedefaultmidpunct}
{\mcitedefaultendpunct}{\mcitedefaultseppunct}\relax
\EndOfBibitem
\bibitem[Chang(2011)]{chang2011}
Chang,~Z. \emph{Fundamentals of Attosecond Optics}, 1st ed.; CRC Press,
  2011\relax
\mciteBstWouldAddEndPuncttrue
\mciteSetBstMidEndSepPunct{\mcitedefaultmidpunct}
{\mcitedefaultendpunct}{\mcitedefaultseppunct}\relax
\EndOfBibitem
\bibitem[Varj\'{u} \latin{et~al.}(2005)Varj\'{u}, Mairesse, Carr\'{e}, Gaarde,
  Johnsson, Kazamias, L\'{o}pez-Martens, Mauritsson, Schafer, Balcou,
  L'Huillier, and Salières]{varju2005}
Varj\'{u},~K.; Mairesse,~Y.; Carr\'{e},~B.; Gaarde,~M.~B.; Johnsson,~P.;
  Kazamias,~S.; L\'{o}pez-Martens,~R.; Mauritsson,~J.; Schafer,~K.~J.;
  Balcou,~P.; L'Huillier,~A.; Salières,~P. Frequency chirp of harmonic and
  attosecond pulses. \emph{Journal of Modern Optics} \textbf{2005}, \emph{52},
  379--394\relax
\mciteBstWouldAddEndPuncttrue
\mciteSetBstMidEndSepPunct{\mcitedefaultmidpunct}
{\mcitedefaultendpunct}{\mcitedefaultseppunct}\relax
\EndOfBibitem
\bibitem[Niikura \latin{et~al.}(2005)Niikura, Villeneuve, and
  Corkum]{niikura2005}
Niikura,~H.; Villeneuve,~D.~M.; Corkum,~P.~B. Mapping Attosecond Electron Wave
  Packet Motion. \emph{Phys. Rev. Lett.} \textbf{2005}, \emph{94}, 083003\relax
\mciteBstWouldAddEndPuncttrue
\mciteSetBstMidEndSepPunct{\mcitedefaultmidpunct}
{\mcitedefaultendpunct}{\mcitedefaultseppunct}\relax
\EndOfBibitem
\bibitem[S\'{a}ndor \latin{et~al.}(2019)S\'{a}ndor, Sissay, Mauger, Gordon,
  Gorman, Scarborough, Gaarde, Lopata, Schafer, and Jones]{sandor2019}
S\'{a}ndor,~P.; Sissay,~A.; Mauger,~F.; Gordon,~M.~W.; Gorman,~T.~T.;
  Scarborough,~T.~D.; Gaarde,~M.~B.; Lopata,~K.; Schafer,~K.~J.; Jones,~R.~R.
  Angle-dependent strong-field ionization of halomethanes. \emph{The Journal of
  Chemical Physics} \textbf{2019}, \emph{151}, 194308\relax
\mciteBstWouldAddEndPuncttrue
\mciteSetBstMidEndSepPunct{\mcitedefaultmidpunct}
{\mcitedefaultendpunct}{\mcitedefaultseppunct}\relax
\EndOfBibitem
\bibitem[foo()]{footnote1}
Preliminary TDDFT calculations in $\text{BrC}_{4}\text{H}$ including Ehrenfest
  nuclear dynamics indicate that the molecule is quite rigid.\relax
\mciteBstWouldAddEndPunctfalse
\mciteSetBstMidEndSepPunct{\mcitedefaultmidpunct}
{}{\mcitedefaultseppunct}\relax
\EndOfBibitem
\bibitem[Perdew and Zunger(1981)Perdew, and Zunger]{perdew1981}
Perdew,~J.~P.; Zunger,~A. Self-interaction correction to density-functional
  approximations for many-electron systems. \emph{Phys. Rev. B} \textbf{1981},
  \emph{23}, 5048--5079\relax
\mciteBstWouldAddEndPuncttrue
\mciteSetBstMidEndSepPunct{\mcitedefaultmidpunct}
{\mcitedefaultendpunct}{\mcitedefaultseppunct}\relax
\EndOfBibitem
\bibitem[Yabana and Bertsch(1996)Yabana, and Bertsch]{yabana1996}
Yabana,~K.; Bertsch,~G.~F. Time-dependent local-density approximation in real
  time. \emph{Phys. Rev. B} \textbf{1996}, \emph{54}, 4484--4487\relax
\mciteBstWouldAddEndPuncttrue
\mciteSetBstMidEndSepPunct{\mcitedefaultmidpunct}
{\mcitedefaultendpunct}{\mcitedefaultseppunct}\relax
\EndOfBibitem
\bibitem[Marques \latin{et~al.}(2012)Marques, Oliveira, and
  Burnus]{marques2012}
Marques,~M.~A.; Oliveira,~M.~J.; Burnus,~T. Libxc: A library of exchange and
  correlation functionals for density functional theory. \emph{Computer Physics
  Communications} \textbf{2012}, \emph{183}, 2272--2281\relax
\mciteBstWouldAddEndPuncttrue
\mciteSetBstMidEndSepPunct{\mcitedefaultmidpunct}
{\mcitedefaultendpunct}{\mcitedefaultseppunct}\relax
\EndOfBibitem
\bibitem[Legrand \latin{et~al.}(2002)Legrand, Suraud, and
  Reinhard]{legrand2002}
Legrand,~C.; Suraud,~E.; Reinhard,~P.-G. Comparison of
  self-interaction-corrections for metal clusters. \emph{Journal of Physics B:
  Atomic, Molecular and Optical Physics} \textbf{2002}, \emph{35},
  1115--1128\relax
\mciteBstWouldAddEndPuncttrue
\mciteSetBstMidEndSepPunct{\mcitedefaultmidpunct}
{\mcitedefaultendpunct}{\mcitedefaultseppunct}\relax
\EndOfBibitem
\bibitem[Ciofini \latin{et~al.}(2003)Ciofini, Chermette, and
  Adamo]{ciofini2003}
Ciofini,~I.; Chermette,~H.; Adamo,~C. A mean-field self-interaction correction
  in density functional theory: implementation and validation for molecules.
  \emph{Chemical Physics Letters} \textbf{2003}, \emph{380}, 12--20\relax
\mciteBstWouldAddEndPuncttrue
\mciteSetBstMidEndSepPunct{\mcitedefaultmidpunct}
{\mcitedefaultendpunct}{\mcitedefaultseppunct}\relax
\EndOfBibitem
\bibitem[Tsuneda and Hirao(2014)Tsuneda, and Hirao]{tsuneda2014}
Tsuneda,~T.; Hirao,~K. Self-interaction corrections in density functional
  theory. \emph{The Journal of Chemical Physics} \textbf{2014}, \emph{140},
  18A513\relax
\mciteBstWouldAddEndPuncttrue
\mciteSetBstMidEndSepPunct{\mcitedefaultmidpunct}
{\mcitedefaultendpunct}{\mcitedefaultseppunct}\relax
\EndOfBibitem
\bibitem[Marques \latin{et~al.}(2003)Marques, Castro, Bertsch, and
  Rubio]{marques2003}
Marques,~M.~A.; Castro,~A.; Bertsch,~G.~F.; Rubio,~A. octopus: a
  first-principles tool for excited electron–ion dynamics. \emph{Computer
  Physics Communications} \textbf{2003}, \emph{151}, 60--78\relax
\mciteBstWouldAddEndPuncttrue
\mciteSetBstMidEndSepPunct{\mcitedefaultmidpunct}
{\mcitedefaultendpunct}{\mcitedefaultseppunct}\relax
\EndOfBibitem
\bibitem[Castro \latin{et~al.}(2006)Castro, Appel, Oliveira, Rozzi, Andrade,
  Lorenzen, Marques, Gross, and Rubio]{castro2006}
Castro,~A.; Appel,~H.; Oliveira,~M.; Rozzi,~C.~A.; Andrade,~X.; Lorenzen,~F.;
  Marques,~M. A.~L.; Gross,~E. K.~U.; Rubio,~A. octopus: a tool for the
  application of time-dependent density functional theory. \emph{physica status
  solidi (b)} \textbf{2006}, \emph{243}, 2465--2488\relax
\mciteBstWouldAddEndPuncttrue
\mciteSetBstMidEndSepPunct{\mcitedefaultmidpunct}
{\mcitedefaultendpunct}{\mcitedefaultseppunct}\relax
\EndOfBibitem
\bibitem[Andrade \latin{et~al.}(2015)Andrade, Strubbe, De~Giovannini, Larsen,
  Oliveira, Alberdi-Rodriguez, Varas, Theophilou, Helbig, Verstraete, Stella,
  Nogueira, Aspuru-Guzik, Castro, Marques, and Rubio]{andrade2015}
Andrade,~X. \latin{et~al.}  Real-space grids and the Octopus code as tools for
  the development of new simulation approaches for electronic systems.
  \emph{Phys. Chem. Chem. Phys.} \textbf{2015}, \emph{17}, 31371--31396\relax
\mciteBstWouldAddEndPuncttrue
\mciteSetBstMidEndSepPunct{\mcitedefaultmidpunct}
{\mcitedefaultendpunct}{\mcitedefaultseppunct}\relax
\EndOfBibitem
\bibitem[Tancogne-Dejean \latin{et~al.}(2020)Tancogne-Dejean, Oliveira,
  Andrade, Appel, Borca, Le~Breton, Buchholz, Castro, Corni, Correa,
  De~Giovannini, Delgado, Eich, Flick, Gil, Gomez, Helbig, Hübener, Jestädt,
  Jornet-Somoza, Larsen, Lebedeva, Lüders, Marques, Ohlmann, Pipolo, Rampp,
  Rozzi, Strubbe, Sato, Schäfer, Theophilou, Welden, and
  Rubio]{tancogne-dejean2020}
Tancogne-Dejean,~N. \latin{et~al.}  Octopus, a computational framework for
  exploring light-driven phenomena and quantum dynamics in extended and finite
  systems. \emph{The Journal of Chemical Physics} \textbf{2020}, \emph{152},
  124119\relax
\mciteBstWouldAddEndPuncttrue
\mciteSetBstMidEndSepPunct{\mcitedefaultmidpunct}
{\mcitedefaultendpunct}{\mcitedefaultseppunct}\relax
\EndOfBibitem
\bibitem[Hamer \latin{et~al.}(2021)Hamer, Tuthill, Scarborough, DiMauro,
  Lopata, Schafer, Gaarde, and Mauger]{hamer2021}
Hamer,~K.~A.; Tuthill,~D.~R.; Scarborough,~T.~D.; DiMauro,~L.~F.; Lopata,~K.;
  Schafer,~K.~J.; Gaarde,~M.~B.; Mauger,~F. Orbital-resolved calculations of
  two-center interferences in linear triatomic molecules. \emph{Phys. Rev. A}
  \textbf{2021}, \emph{104}, 033114\relax
\mciteBstWouldAddEndPuncttrue
\mciteSetBstMidEndSepPunct{\mcitedefaultmidpunct}
{\mcitedefaultendpunct}{\mcitedefaultseppunct}\relax
\EndOfBibitem
\end{mcitethebibliography}

\newpage
\section{TOC Image}
\hfill

\begin{figure}[!h]
    \centering
    \includegraphics[width=0.9\linewidth]{Figures/toc.png}
    \label{fig:toc}
\end{figure}

\end{document}